\documentclass[11pt]{article}
% Preambulo comun para los informes de la ronda 3.
% Se incluye con % Preambulo comun para los informes de la ronda 3.
% Se incluye con % Preambulo comun para los informes de la ronda 3.
% Se incluye con \input{../comun/preambulo} desde cada subcarpeta.
\usepackage[utf8]{inputenc}
\usepackage[T1]{fontenc}
\usepackage[spanish,es-nodecimaldot]{babel}
\usepackage{lmodern}
\usepackage[a4paper,margin=2.4cm]{geometry}
\usepackage{booktabs}
\usepackage{longtable}
\usepackage{array}
\usepackage{xcolor}
\usepackage{hyperref}
\usepackage{enumitem}
\usepackage{fancyhdr}
\usepackage{titlesec}
\usepackage{graphicx}
\usepackage{amssymb}
\usepackage{microtype}

\definecolor{sosTeal}{HTML}{0F766E}
\definecolor{sosGray}{HTML}{475569}

\hypersetup{
  colorlinks=true,
  linkcolor=sosTeal,
  urlcolor=sosTeal,
  citecolor=sosTeal,
  pdftitle={SOS Clinic - Ronda 3 de pruebas},
  pdfauthor={Equipo SOS Clinic}
}

\pagestyle{fancy}
\fancyhf{}
\lhead{\textcolor{sosGray}{SOS Clinic --- Ronda 3}}
\rhead{\textcolor{sosGray}{\leftmark}}
\cfoot{\thepage}
\setlength{\headheight}{15pt}
\renewcommand{\headrulewidth}{0.4pt}

\titleformat{\section}{\Large\bfseries\color{sosTeal}}{\thesection}{0.8em}{}
\titleformat{\subsection}{\large\bfseries\color{sosGray}}{\thesubsection}{0.6em}{}

\setlist[itemize]{leftmargin=1.4em,itemsep=2pt,topsep=4pt}
\setlist[enumerate]{leftmargin=1.6em,itemsep=2pt,topsep=4pt}

\newcommand{\pantalla}[1]{\texttt{\detokenize{#1}}}
\newcommand{\ok}{\textcolor{sosTeal}{\textbf{PASS}}}
\newcommand{\fallo}{\textcolor{red!70!black}{\textbf{FAIL}}}
\newcommand{\aviso}{\textcolor{orange!80!black}{\textbf{WARN}}}
\newcommand{\omitido}{\textcolor{sosGray}{\textbf{SKIP}}}
 desde cada subcarpeta.
\usepackage[utf8]{inputenc}
\usepackage[T1]{fontenc}
\usepackage[spanish,es-nodecimaldot]{babel}
\usepackage{lmodern}
\usepackage[a4paper,margin=2.4cm]{geometry}
\usepackage{booktabs}
\usepackage{longtable}
\usepackage{array}
\usepackage{xcolor}
\usepackage{hyperref}
\usepackage{enumitem}
\usepackage{fancyhdr}
\usepackage{titlesec}
\usepackage{graphicx}
\usepackage{amssymb}
\usepackage{microtype}

\definecolor{sosTeal}{HTML}{0F766E}
\definecolor{sosGray}{HTML}{475569}

\hypersetup{
  colorlinks=true,
  linkcolor=sosTeal,
  urlcolor=sosTeal,
  citecolor=sosTeal,
  pdftitle={SOS Clinic - Ronda 3 de pruebas},
  pdfauthor={Equipo SOS Clinic}
}

\pagestyle{fancy}
\fancyhf{}
\lhead{\textcolor{sosGray}{SOS Clinic --- Ronda 3}}
\rhead{\textcolor{sosGray}{\leftmark}}
\cfoot{\thepage}
\setlength{\headheight}{15pt}
\renewcommand{\headrulewidth}{0.4pt}

\titleformat{\section}{\Large\bfseries\color{sosTeal}}{\thesection}{0.8em}{}
\titleformat{\subsection}{\large\bfseries\color{sosGray}}{\thesubsection}{0.6em}{}

\setlist[itemize]{leftmargin=1.4em,itemsep=2pt,topsep=4pt}
\setlist[enumerate]{leftmargin=1.6em,itemsep=2pt,topsep=4pt}

\newcommand{\pantalla}[1]{\texttt{\detokenize{#1}}}
\newcommand{\ok}{\textcolor{sosTeal}{\textbf{PASS}}}
\newcommand{\fallo}{\textcolor{red!70!black}{\textbf{FAIL}}}
\newcommand{\aviso}{\textcolor{orange!80!black}{\textbf{WARN}}}
\newcommand{\omitido}{\textcolor{sosGray}{\textbf{SKIP}}}
 desde cada subcarpeta.
\usepackage[utf8]{inputenc}
\usepackage[T1]{fontenc}
\usepackage[spanish,es-nodecimaldot]{babel}
\usepackage{lmodern}
\usepackage[a4paper,margin=2.4cm]{geometry}
\usepackage{booktabs}
\usepackage{longtable}
\usepackage{array}
\usepackage{xcolor}
\usepackage{hyperref}
\usepackage{enumitem}
\usepackage{fancyhdr}
\usepackage{titlesec}
\usepackage{graphicx}
\usepackage{amssymb}
\usepackage{microtype}

\definecolor{sosTeal}{HTML}{0F766E}
\definecolor{sosGray}{HTML}{475569}

\hypersetup{
  colorlinks=true,
  linkcolor=sosTeal,
  urlcolor=sosTeal,
  citecolor=sosTeal,
  pdftitle={SOS Clinic - Ronda 3 de pruebas},
  pdfauthor={Equipo SOS Clinic}
}

\pagestyle{fancy}
\fancyhf{}
\lhead{\textcolor{sosGray}{SOS Clinic --- Ronda 3}}
\rhead{\textcolor{sosGray}{\leftmark}}
\cfoot{\thepage}
\setlength{\headheight}{15pt}
\renewcommand{\headrulewidth}{0.4pt}

\titleformat{\section}{\Large\bfseries\color{sosTeal}}{\thesection}{0.8em}{}
\titleformat{\subsection}{\large\bfseries\color{sosGray}}{\thesubsection}{0.6em}{}

\setlist[itemize]{leftmargin=1.4em,itemsep=2pt,topsep=4pt}
\setlist[enumerate]{leftmargin=1.6em,itemsep=2pt,topsep=4pt}

\newcommand{\pantalla}[1]{\texttt{\detokenize{#1}}}
\newcommand{\ok}{\textcolor{sosTeal}{\textbf{PASS}}}
\newcommand{\fallo}{\textcolor{red!70!black}{\textbf{FAIL}}}
\newcommand{\aviso}{\textcolor{orange!80!black}{\textbf{WARN}}}
\newcommand{\omitido}{\textcolor{sosGray}{\textbf{SKIP}}}


\title{\textbf{Análisis estático y de dependencias}\\[0.3em]
\large Metodología SAST/SCA --- Ronda 3}
\author{Equipo SOS Clinic}
\date{22 de octubre de 2026}

\begin{document}
\maketitle

\section{Objetivo}

Buscar secretos versionados, patrones inseguros en el código y vulnerabilidades
conocidas en las dependencias, sin modificar el repositorio y sin enviar código
a servicios externos más allá de la descarga de reglas y bases de datos de las
herramientas.

\section{Herramientas}

\begin{table}[h]
\centering
\small
\begin{tabular}{@{}p{2.6cm}p{4.2cm}p{7.2cm}@{}}
\toprule
Herramienta & Imagen / versión & Uso \\
\midrule
Gitleaks & \texttt{ghcr.io/gitleaks/gitleaks:latest} & Detección de secretos en el árbol y en el historial. \\
Semgrep & \texttt{semgrep/semgrep:latest} & Reglas \texttt{p/security-audit} y \texttt{p/secrets}. \\
Trivy & \texttt{aquasec/trivy:latest} & Vulnerabilidades, secretos y configuración. \\
\bottomrule
\end{tabular}
\caption{Herramientas ejecutadas vía Docker el 22 de octubre de 2026.}
\end{table}

\section{Alcance}

\begin{itemize}
  \item Árbol versionado: 302 archivos, aproximadamente 15,7 MB.
  \item Historial de Git: 45 commits.
  \item Dependencias: los 7 archivos \texttt{pom.xml}.
  \item Infraestructura: \texttt{infra/docker-compose.yml} y configuración de
  los servicios.
\end{itemize}

\section{Procedimiento}

\subsection{Gitleaks}

\begin{enumerate}
  \item Escaneo del árbol con \texttt{--no-git}, la configuración
  \texttt{.gitleaks.toml} del repositorio y \texttt{--redact} para no imprimir
  valores sensibles.
  \item Escaneo del historial completo con la misma configuración y redacción.
  \item Revisión de las reglas y archivos reportados para distinguir hallazgos
  vigentes de históricos.
\end{enumerate}

\subsection{Semgrep}

\begin{enumerate}
  \item Ejecución con las reglas \texttt{p/security-audit} y \texttt{p/secrets},
  \texttt{--metrics off} y salida JSON.
  \item Clasificación manual de los hallazgos por archivo y contexto: código de
  producción, scripts de prueba o configuración.
  \item Verificación de que los hashes BCrypt reportados son fixtures de pruebas
  y no credenciales reales.
\end{enumerate}

\subsection{Trivy}

\begin{enumerate}
  \item Escaneo \texttt{fs} con los escáneres \texttt{vuln}, \texttt{secret} y
  \texttt{misconfig}.
  \item La primera ejecución falló con HTTP 429 de Maven Central al resolver
  los BOM. Se repitió montando la caché local \texttt{\textasciitilde/.m2} y
  con \texttt{--offline-scan}.
  \item Escaneo específico de configuración sobre \texttt{infra/} y revisión
  manual del compose ante la ausencia de archivos detectados.
\end{enumerate}

\section{Criterios}

\begin{itemize}
  \item \ok{} si no hay hallazgos vigentes o si los hallazgos son falsos
  positivos demostrables.
  \item \aviso{} si el hallazgo es histórico, está en un archivo eliminado o
  requiere una decisión de rotación.
  \item \fallo{} si hay un secreto vigente, una dependencia crítica/alta sin
  justificar o un patrón inseguro en código de producción.
\end{itemize}

\section{Limitaciones}

\begin{itemize}
  \item Las imágenes usan la etiqueta \texttt{latest}; los resultados
  corresponden a las versiones descargadas el 22 de octubre de 2026.
  \item Las reglas comunitarias de Semgrep no cubren todos los patrones
  posibles; no sustituyen una revisión manual.
  \item Trivy no detectó \texttt{docker-compose.yml} como archivo de
  configuración; se revisó a mano.
  \item El escaneo de secretos del historial puede marcar valores ya rotados o
  archivos eliminados; se clasifican caso por caso.
\end{itemize}

\section{Reproducción}

\begin{verbatim}
docker run --rm -v <repo>:/repo ghcr.io/gitleaks/gitleaks:latest \
  detect --source /repo --config /repo/.gitleaks.toml --redact --no-git

docker run --rm -v <repo>:/src semgrep/semgrep \
  semgrep scan --config p/security-audit --config p/secrets \
  --metrics off /src

docker run --rm -v <repo>:/src -v ~/.m2:/root/.m2:ro \
  aquasec/trivy fs --scanners vuln,secret,misconfig \
  --offline-scan /src
\end{verbatim}

\end{document}
